\documentclass[11pt]{article}
\usepackage[T1]{fontenc}
\usepackage[utf8]{inputenc}
\usepackage{lmodern,microtype,geometry,amsmath,amssymb,amsthm,mathtools,enumitem}
\usepackage[colorlinks=true,linkcolor=blue,urlcolor=blue,citecolor=blue]{hyperref}
\hypersetup{pdftitle={Quadratic-width learning with symmetry-invariant neuron merging},pdfauthor={ChatGPT-6 Astra; Claude Fable 5 and Opus 5; Kaarel Hanni; Dmitry Vaintrob}}
\geometry{margin=1in}
\newtheorem{theorem}{Theorem}
\newtheorem{lemma}[theorem]{Lemma}
\newtheorem{proposition}[theorem]{Proposition}
\newtheorem{corollary}[theorem]{Corollary}
\newtheorem{remark}[theorem]{Remark}
\newcommand{\R}{\mathbb R}
\newcommand{\E}{\mathbb E}
\newcommand{\Prb}{\mathbb P}
\newcommand{\norm}[1]{\left\|#1\right\|}
\newcommand{\clip}{\operatorname{clip}}
\title{Quadratic-width learning with\protect\\symmetry-invariant neuron merging}
\author{ChatGPT-6 Astra \and Claude Fable 5 and Opus 5 \and Kaarel H\"anni \and Dmitry Vaintrob\thanks{Authors are listed in alphabetical order.}}
\date{Research draft, 21 October 2026}
\begin{document}
\maketitle
\begin{abstract}
We strengthen a quadratic-width Bayesian learning theorem by removing two
artifacts of its neuron-merging restriction: dependence on arbitrary hidden
neuron signs and the requirement of nearly equal population sizes. The
new condition allows signed populations of arbitrary sizes, with different
population counts in different layers. A mass-normalized representation
makes every compressed Gaussian coordinate have the same variance,
independently of the smallest population. A scale-normalized readout budget
controls the number of populations. The resulting sample bound remains
quadratic in Boolean teacher width, up to logarithms, with polynomial depth
dependence, at every learner width and every specified Gaussian scale at or
above mean field. We also give a local operator-norm certificate: approximate
agreement of rows on population-constant states, together with internal-noise
robustness, implies prediction-preserving merging. These are structural
prior restrictions; the paper does not establish that ordinary trained
networks satisfy them. The results are not Lean verified.
\end{abstract}

\section{Which aspects of neuron merging are natural?}

Clustering redundant neurons is a plausible way to reduce a large network
to a smaller representation. A theorem based on this idea should respect
exact symmetries of the network, and should not require every cluster to
have almost the same size. Both requirements matter mathematically.

The previous population-averaging condition can distinguish two identical
functions solely because one hidden neuron has been sign-flipped. Take
$d=1$, $N=2$, one hidden layer, first-layer rows $(2,0),(2,0)$, and readout
$(3/8,3/8)$. The output is $(3/4)\tanh(2x)$ and $E_v=N\norm v_2^2=9/16$.
The previous budget $q\le1+E_v$ permits only one population, whose averaging
preserves this network exactly. Flip the second row and its readout sign.
The function and Gaussian density are unchanged, but one-population
averaging now gives the zero function, changing clipped scores by $1/2$.
Both representations satisfy the internal-noise condition below at $c=0.1$:
the calibrated ball changes their raw outputs by at most $0.1$, and the
clean output magnitude is greater than $0.7$.

Allowing signed populations repairs this defect. Allowing arbitrary group
sizes is more subtle: ordinary compressed weights can become large when a
population is small. We remove this difficulty by rescaling compressed
hidden states by the square root of their population mass. The resulting
activation family is uniformly Lipschitz in its scale parameter, so even
very small populations admit a finite code independent of learner width.

The final section gives a local sufficient condition for merging. It uses
operator norms of explicit weight defects instead of assuming directly that
predictions are preserved. The empirical prevalence of either condition
remains an open question.

\section{Model and a symmetry-respecting support condition}

Let $H=L-1\ge1$, $d\ge1$, and $\mathcal X=\{-1,1\}^d$. The learner has
$H$ hidden tanh layers of width $N$, first-layer augmented input
$\bar x=(x,1)$, and linear readout:
\[
 h^1=\tanh(W^1\bar x),\qquad h^s=\tanh(W^sh^{s-1}),\qquad
 z_w(x)=v^\top h^H.
\]
Under $\mu_{N,t}$, independent Gaussian entries have variances
$t^2/d$, $t^2/N$, and $t^2/N^2$ in the first layer, subsequent hidden
layers, and readout, respectively; $0<t\le1$. Taking
$t=N^{(1-\gamma)/2}$ includes every $\gamma\ge1$. Set
\[
 E_v=N\norm v_2^2,\qquad M(w)=\max_{x\in\mathcal X}|z_w(x)|,
 \qquad T(u)=\clip(u,-1/2,1/2).
\]

At hidden layer $s$, choose a partition into $q_s$ nonempty groups of
arbitrary sizes $k_{s,a}$ and a sign $\sigma_i\in\{-1,1\}$ for each neuron.
Let $F_s$ have column $a$ equal to the signed normalized indicator of that
population: its entries there are $\sigma_i/\sqrt{k_{s,a}}$. Then
$F_s^\top F_s=I$, and $R_s=F_sF_s^\top$ averages neurons after aligning
their signs. Define the merged network by
\begin{equation}\label{eq:projection}
 \bar W^1=R_1W^1,\qquad
 \bar W^s=R_sW^sR_{s-1},\qquad \bar v=R_Hv.
\end{equation}
Write $q=\max_s q_s$. The new event $\mathcal A_N$ requires the existence
of such signed populations satisfying
\begin{equation}\label{eq:population-condition}
 (q-1)M(w)^2\le4E_v(w),\qquad
 \max_x|T(z_w(x))-T(z_{\bar w}(x))|\le1/16.
\end{equation}
The division-free budget includes $M=0$ without a special convention.
It is invariant under readout rescaling. The clipped-score comparison uses
fixed margins and retains their output-scale dependence. The complete
condition is invariant under hidden-neuron sign flips and permutations.
Indeed, a sign change $S_s$ replaces $F_s$ by $S_sF_s$ and preserves all
functions and energies in the condition. All witnesses are finite choices
at fixed $N$, so the event is Borel.

We can additionally impose internal-noise robustness. Add disturbances
after every activation and at the output, with
\[
 \norm\eta_*^2=\sum_s\norm{\eta^s}_2^2/N+|\eta^{\rm out}|^2,
 \qquad r_c(w)=c/\sqrt{1+HE_v(w)},\quad c>0.
\]
Let $\mathcal R_{N,c}$ require
\begin{equation}\label{eq:rms-condition}
 |T(z_w(x;\eta))-T(z_w(x))|\le1/16
 \quad(x\in\mathcal X,\ \norm\eta_*\le r_c(w)).
\end{equation}
The prior is the raw Gaussian conditioned on $\mathcal A_N$, optionally
also on $\mathcal R_{N,c}$. Training further conditions on
$f(X_i)z_w(X_i)>1/2$. All statements concern exact Bayesian conditioning.

\section{An exact normalized representation}

Put $a_{s,b}=\sqrt{k_{s,b}/N}$ and define
\[
 B^1=F_1^\top W^1/\sqrt N,\qquad
 B^s=F_s^\top W^sF_{s-1},\qquad b=\sqrt N F_H^\top v.
\]
The merged network has exactly the following compressed representation:
\begin{equation}\label{eq:perspective-network}
 g^1=\phi_{a_1}(B^1\bar x),\qquad
 g^s=\phi_{a_s}(B^sg^{s-1}),\qquad z_{\bar w}=b^\top g^H,
\end{equation}
where the activation is applied coordinatewise and
\[
 \phi_a(u)=a\tanh(u/a)\quad(a>0),\qquad \phi_0(u)=0.
\]
To verify the identity, the full hidden state is $\sqrt N F_sg^s$;
oddness of tanh gives \eqref{eq:perspective-network} at every layer.

\begin{lemma}[Lipschitz scale coordinates]\label{lem:perspective}
For $a,b\ge0$ and real $u,v$,
\[
 |\phi_a(u)-\phi_a(v)|\le|u-v|,\qquad
 |\phi_a(u)-\phi_b(u)|\le|a-b|.
\]
\end{lemma}
\begin{proof}
The input derivative is $\operatorname{sech}^2(u/a)$. The scale derivative
is $\tanh y-y\operatorname{sech}^2y$, where $y=u/a$. It is odd and,
for $y\ge0$, increases from 0 to 1: its derivative is
$2y\operatorname{sech}^2y\tanh y\ge0$. Continuity at $a=0$ follows from
$|\phi_a(u)|\le a$.
\end{proof}

For fixed populations, every entry of $\sqrt d B^1$, $B^s$ for $s\ge2$,
and $b$ is an independent Gaussian of variance exactly $t^2/N$.
This follows from the orthonormal columns of $F_s$, and is independent of
the group sizes. Write
\[
 \kappa=d\norm{B^1}_F^2+\sum_{s=2}^H\norm{B^s}_F^2+\norm b_2^2,
 \quad P=q_1(d+1)+\sum_{s=2}^Hq_sq_{s-1}+q_H.
\]
Then
\begin{equation}\label{eq:mgf}
 \E e^{N\kappa/(4t^2)}=2^{P/2},\qquad
 P/N\le d+2+(H-1)Q\quad(q_s\le\min(N,Q)).
\end{equation}
The auxiliary representation uses perspective activations; the learner
itself continues to use ordinary tanh.

\section{A quadratic sample bound}

The Boolean teacher has width $C$ and depth $H$, with NAND, copy, and
negation gates. Constant wires and basis-compilation overhead count toward
these resources. Define
\begin{align}
 e_*&=(16\sqrt{2C})^{-1},&\tau_0&=e_*/4,\notag\\
 b_L&=\sqrt{16\log9+8\log(4L)},&A&=\max\{32b_L,2\log(8/e_*)\},\notag\\
 a&=\max\{4,256c\},&\rho&=\min\{1/4,A/[4(d+1)],Ae_*/(64\sqrt C)\},\notag\\
 q_0&=d+2+(H-1)C,&B_0&=9dA^2+18(H-1)CA^2+2a^2C,\notag\\
 D_*&=B_0+q_0\log\frac{e\sqrt{2\pi}}{2\rho}+\log2.
 \label{eq:target-constants}
\end{align}
Appendix~\ref{app:cylinder} proves that the target has a cylinder of raw
probability at least $e^{-ND_*/t^2}$ in the stated support, for every
$N\ge C$. In particular
$D_*=O((d+LC)\log^2(2dLC)+c^2C)$.

For $0<\varepsilon,\delta<1$, set
\begin{align}
 B&=\max\{8,8[D_*+\log(4/\varepsilon)]\},\qquad Q=\lceil1+16B\rceil,\notag\\
 J&=4\left[D_*+\log(4/\varepsilon)+H\log Q+H\log(2Q)
       +\tfrac{\log2}{2}\{d+2+(H-1)Q\}\right],\notag\\
 R&=1+\sqrt J,\qquad A_g=2+d+(Q+1)(R+1),\notag\\
 \zeta&=[32(H+1)A_g^{H+1}]^{-1},\qquad
 D_Q=Q(d+2)+(H-1)Q^2+HQ.\label{eq:budget}
\end{align}

\begin{theorem}\label{thm:main}
If
\[
 n\ge\frac2\varepsilon
       [D_Q\log(3+4R/\zeta)+\log(1/\delta)],
\]
then, with probability at least $1-\delta$ over iid inputs from any
fixed distribution $D$, the posterior satisfies
\[
 \E_{w\sim Q_{S,N,t,c}}\Prb_{X\sim D}\{f(X)z_w(X)\le1/4\}\le\varepsilon
\]
simultaneously for every $N\ge C$ and $0<t\le1$. All normalizers are
positive. The theorem applies with or without the additional condition
$\mathcal R_{N,c}$.
\end{theorem}
\begin{proof}
The target cylinder lies in every noiseless training event, so the posterior
normalizer relative to the raw Gaussian is at least $e^{-ND_*/t^2}$.
Since $NE_v/t^2$ is chi-square with $N$ degrees of freedom,
\[
 \mu_{N,t}(E_v>B)\le e^{-NB/(8t^2)}\quad(B\ge8).
\]
After normalization this event has probability at most $\varepsilon/4$.
A nonempty successful training sample forces $M>1/2$; therefore
\eqref{eq:population-condition} implies $q<1+16E_v$. On the retained
posterior, all witnessing population counts are at most $Q$.

There are at most $Q^H(2Q)^{HN}$ tuples of signed partitions with each
$q_s\le\min(N,Q)$. From \eqref{eq:mgf}, the raw probability that any such
tuple has $\kappa>J$ is at most
\[
 \exp\left(H\log Q+HN\log(2Q)
 +\tfrac N2[d+2+(H-1)Q]\log2-\frac{NJ}{4t^2}\right).
\]
The choice of $J$, using $N/t^2\ge1$ and $t\le1$, bounds this by
$e^{-ND_*/t^2}e^{-(N/t^2)\log(4/\varepsilon)}$. Its posterior probability
is at most $\varepsilon/4$. This is a union over all partitions before
selecting a witness from the weights.

Every retained compressed coefficient is bounded by $\sqrt J$, and every
scale is in $[0,1]$. Pad to width $Q$ with zero scales. Round coefficients
and scales to fixed grids with error at most $\zeta$, preserving zero and
the interval $[0,1]$ for scales. Lemma~\ref{lem:perspective} gives hidden
coordinate errors
\[
 e_1\le(d+2)\zeta,\qquad
 e_s\le QR e_{s-1}+(Q+1)\zeta.
\]
The output error is at most $QR e_H+Q\zeta\le1/32$. These rounded
functions form a class $\mathcal G$, independent of $N,t$, with
\[
 \log|\mathcal G|\le D_Q\log(3+4R/\zeta).
\]
The score difference from the original predictor is at most
$1/16+1/32=3/32$.

Use a ramp loss on $f(x)T(h(x))$ equal to one below $11/32$, zero above
$13/32$, and affine between these thresholds. Each retained predictor has
a candidate in $\mathcal G$ with zero empirical ramp loss whose population
ramp loss bounds the original predictor's test-margin error. A finite-class
zero-loss bound gives, with probability at least $1-\delta$,
\[
 \Prb_D\{fz_w\le1/4\}
 \le [\log|\mathcal G|+\log(1/\delta)]/n
\]
for every retained predictor. For completeness, a fixed $[0,1]$-valued loss
of mean $r$ is zero on all $n$ examples with probability at most
$(1-r)^n\le e^{-nr}$; then union bound. Add the discarded posterior mass
$\varepsilon/2$. The candidate class, and hence this sample event, is
independent of learner width and Gaussian scale.
\end{proof}

The sample rate is
\[
 \widetilde O\left(\frac{LdQ+L^2Q^2+\log(1/\delta)}\varepsilon\right),
 \qquad Q=\widetilde O(d+LC+c^2C+\log(1/\varepsilon)).
\]
For a single-output bounded-fan-in teacher, pruning the output's ancestors
gives width $C'\le\min(C,2^{H-1}+1)$, including at most one constant wire
per stage. Thus $\log C'=O(L)$. To make the polynomial degree explicit, put
$S=2+d+L+c+\log(4/\varepsilon)$. The cylinder parameters satisfy
$A,a=O(S)$ and $\log(1/\rho)=O(S)$; hence
\[
 D_*=O(C'S^3),\quad Q=O(C'S^3),\quad J=O(C'S^4),\quad
 D_Q=O(C'^2S^7),
\]
while $\log(3+4R/\zeta)=O(S^2)$. Consequently there is an absolute
constant $K_0$ for which
\[
 n\ge\frac{K_0}\varepsilon
 [C^2(2+d+L+c+\log(4/\varepsilon))^9+\log(1/\delta)]
\]
is sufficient. The extra $HQ$ scale parameters do not change this degree.

\section{A local weight certificate for merging}

Signed population projections commute with the property of being a tanh
hidden state: applying tanh to a vector in their range keeps it in their
range. For proposed signed projections $R_s$, define
\[
 \Delta_1=\sqrt{(d+1)/N}\norm{(I-R_1)W^1}_{\rm op},\qquad
 \Delta_s=\norm{(I-R_s)W^sR_{s-1}}_{\rm op}\quad(s\ge2).
\]
These measure disagreement of rows after restricting their inputs to
population-constant states. They permit arbitrary weights on directions
orthogonal to the preceding populations.

\begin{proposition}[Local agreement implies prediction stability]
If $w\in\mathcal R_{N,c}$ and $\sum_s\Delta_s^2\le r_c(w)^2$, then the
merged network obeys the score comparison in
\eqref{eq:population-condition}.
\end{proposition}
\begin{proof}
Run the merged network. Its preceding hidden state lies in the range of
$R_{s-1}$ and has norm at most $\sqrt N$. The difference between its next
hidden state and $\tanh(W^s\bar h^{s-1})$ has RMS norm at most $\Delta_s$,
by tanh's Lipschitz property. The first-layer bound follows from
$\norm{\bar x}_2=\sqrt{d+1}$. Thus the complete merged trajectory is the
original network under postactivation disturbances of total norm at most
$(\sum_s\Delta_s^2)^{1/2}$. The readout agrees because the last hidden
state lies in the range of $R_H$. Apply \eqref{eq:rms-condition}.
\end{proof}

Consequently we may replace the direct output-comparison support condition
by the existence of populations satisfying the same count budget and this
local certificate, together with $\mathcal R_{N,c}$. The target cylinder
still has enough mass. In \eqref{eq:target-constants}, shrink $\rho$ further
to
\begin{equation}\label{eq:local-radius}
 \rho\le\frac{c}{2\sqrt{1+4Ha^2C}}
 \left[\frac{(d+1)^2}{d}+(H-1)C\right]^{-1/2}.
\end{equation}
On the cylinder, $E_v\le4a^2C$,
$\Delta_1\le\rho(d+1)/\sqrt d$, and
$\Delta_s\le\rho\sqrt C$ for $s\ge2$. Indeed, the teacher is exactly
population-constant, and the Gaussian bulk annihilates the previous
population subspace. Equation~\eqref{eq:local-radius} therefore gives the
local certificate uniformly throughout the cylinder.

Recompute $D_*$ with this smaller radius. Theorem~\ref{thm:main} applies
unchanged. The additional mass cost is logarithmic in the reciprocal radius:
\[
 D_*=O\left((d+LC)[\log^2(2dLC)+\log(2+c+c^{-1})]+c^2C\right).
\]
Thus width dependence remains quadratic up to logarithms. A literal degree-nine
envelope also holds after adding $c^{-1}$ to the displayed parameter sum.
This condition is a concrete approximate agreement of block row sums,
sometimes called an approximate equitable partition or lumpability
condition. It is still restrictive, and no claim about its prevalence in
trained networks follows from the theorem.

\section{Other calibrations and remaining questions}

A purely parameter-based alternative replaces $M(w)$ by $\norm v_1$:
\[
 (q-1)\norm v_1^2\le4N\norm v_2^2.
\]
The same theorem and constants apply. The target cylinder has
$\norm v_1\le a+\rho$, and any successful nonempty training sample forces
$\norm v_1>1/2$. This condition is stronger than the $M$-normalized one.
It uses the standard effective support size
$\norm v_1^2/\norm v_2^2$: the population budget is bounded by four times
the inverse effective fraction of readout neurons, plus one. It can reject
diffuse dense readouts, so easier evaluation should not be confused with
empirical naturality.

The original unnormalized count budget $q\le1+E_v$ is also compatible with
signed arbitrary-size populations: replace $Q$ by $\lceil1+B\rceil$ and
leave the remaining proof unchanged.

\paragraph{Why ordinary dropout needs a separate argument.}
Stability under neuron removal is another natural candidate, but the present
target cylinder does not automatically establish it at mean-field scale.
For example, let a Gaussian bulk row satisfy
\[
 E\sim\mathcal N\!\left(0,\frac{t^2}{N}
       \left[I-\frac{\boldsymbol1\boldsymbol1^\top}{N}\right]\right).
\]
Its clean action on the constant state is exactly zero. If $m$ uniform
coordinates are sampled with replacement and rescaled by $N/m$, the
resulting random state $a$ instead has
\[
 \E(Ea)^2=t^2(N-1)/m.
\]
Indeed, its multinomial count vector gives
$\E\norm a_2^2=N+N(N-1)/m$ and
$\boldsymbol1^\top a=N$. Sampling without replacement gives
$t^2(N/m-1)$. Thus a fixed retained width can amplify Gaussian bulk noise
as $N$ grows. This obstructs a direct transfer of the existing cylinder
proof to dropout stability; it is not a theorem ruling out every dropout
prior or learning guarantee.

These extensions remove exact representation artifacts. They do not prove
that an arbitrary robust network has a small population description. Small
activation distances on a training sample do not by themselves imply the
uniform condition, and arbitrary soft projections need not preserve tanh's
nonlinear structure. The local certificate separates one sufficient weight
condition from these unresolved empirical and structural questions.

\paragraph{Acknowledgments and verification.}
This note builds on the earlier collaborative robust-prior construction,
including contributions from Claude. Independent AI-agent checks examined
the sign counterexample, mass normalization, Gaussian partition bound,
scale-sensitive grid, and local certificate. This is a research draft,
not a Lean-verified theorem or a claim of empirical validation.

\appendix
\section{A target implementation of sufficient prior mass}\label{app:cylinder}

\begin{lemma}[Stable Boolean cylinder]\label{lem:cylinder}
For every $N\ge C$ and $0<t\le1$, there is a set of weights contained in
$\mathcal A_N\cap\mathcal R_{N,c}$ on which $f(x)z_w(x)>1$ for every
input, with raw Gaussian probability at least $\exp(-ND_*/t^2)$.
\end{lemma}
\begin{proof}
Partition the neurons into $C$ balanced clone groups of sizes $k_b$.
Then $C/2\le N/k_b\le2C$. Fill unused logical slots with constants or
copies, so the ideal Boolean state $b^s(x)$ has entries in $\{-1,1\}$,
is constant on each group, and has norm $\sqrt N$. Write $P$ for averaging
within the clone groups.

Use gain $A$ in every gate. NAND has signed affine expression $1-u-v$;
copies and negations use $u$ and $-u$. Count constant wires in internal
layers. Each ideal signed preactivation is at least $A$. Divide each
coefficient from a logical wire equally among its clones, and let the
readout average the output-wire clones with total coefficient $a$.
Denote these teacher weights by $w_*$.

In coordinates $\sqrt d W^1$, $\sqrt N W^s$, and $Nv$, all raw entries
have variance $t^2$. Restrict every first-layer coordinate, every internal
row's coordinate along each normalized clone-group direction, and every
readout coordinate to within $\rho t$ of its teacher value. Internal
matrices then have the form
\[
 W^s=W_*^s+R^s+E^s,\qquad
 R^s=R^sP,\quad \norm{R^s}_{\rm op}\le\rho t\sqrt C,
 \quad E^sP=0.
\]
The unconstrained $E^s$ are independent Gaussian bulk matrices. Two
$1/4$-nets, each of cardinality at most $9^N$, give
\[
 \Prb\{\norm{E^s}_{\rm op}>b_Lt\}
 \le2\exp(2N\log9-Nb_L^2/8)\le\frac1{2L}.
\]
Indeed, each fixed bilinear form has variance at most $t^2/N$, and the
operator norm is at most twice its maximum over the two nets.
All bulk norm bounds therefore hold with conditional probability at least
$1/2$.

We check finite disturbance stability, allowing RMS disturbance $\tau_0$
separately at each hidden layer and scalar disturbance at most $\tau_0$
at the output. This contains the total-budget ball $\norm\eta_*\le\tau_0$.
Put $e_s=\norm{h^s_\eta-b^s}_2/\sqrt N$. If $e_{s-1}\le e_*$, each
group mean differs from its ideal bit by at most
$e_*\sqrt{N/k_b}\le1/16$. At most three teacher coefficients enter a
gate, so its intended signed preactivation remains at least $A/2$.
For $\xi\in\{-1,1\}$ the elementary bound
\[
 \xi q\ge A/2\quad\Longrightarrow\quad
 |\tanh(q+u)-\xi|\le2e^{-A/2}+8|u|/A
\]
follows by considering $|u|<A/4$ and its complement. Since
$E^sb^{s-1}=0$ and $\norm{h^{s-1}_\eta}_2/\sqrt N\le1+\tau_0$,
\[
 e_s\le2e^{-A/2}+(8/A)\rho\sqrt C(1+\tau_0)
                 +(8b_L/A)e_{s-1}+\tau_0\le e_*.
\]
Each term is at most $e_*/4$. In the first layer, the preactivation
perturbation is at most $\rho(d+1)\le A/4$, which initializes the same
invariant. At the readout,
\[
 f(x)z_w(x;\eta)\ge15a/16-\rho(1+\tau_0)-\tau_0>1.
\]
Furthermore, $\sqrt N\norm{v_*}_2\ge a\sqrt{C/2}$ and
$\sqrt N\norm{v-v_*}_2\le\rho$, so
\[
 E_v\ge a^2C/8\ge2C,\qquad
 r_c(w)\le\frac{2\sqrt2c}{a\sqrt C}\le\tau_0.
\]
Thus all noisy and clean clipped outputs on the calibrated ball equal
$f(x)/2$, establishing membership in $\mathcal R_{N,c}$.

The same clone projections in \eqref{eq:projection} map this cylinder
back into itself. Averaging rows preserves the coordinate intervals because
the teacher rows agree within each group. Double projection kills the
bulk, since $E^sP=0$. The first-layer and readout intervals are likewise
preserved. The averaged network therefore also has signed margin greater
than 1 on every input. Its clipped predictions agree exactly with the
original network, and its population count is $C$. Moreover,
\[
 \sqrt{E_v}\ge a\sqrt{C/2}-\rho\ge(5/8)a\sqrt C,\qquad
 M(w)\le\norm v_1\le a+\rho\le5a/4.
\]
Thus $(C-1)M(w)^2\le4E_v$, proving membership in $\mathcal A_N$.

There are $Nq_0$ independent constrained scalar coordinates; the sum of
squares of their centers is at most $NB_0$. For $Z\sim\mathcal N(0,t^2)$,
$\rho\le1/4$, and any center $m$,
\[
 \Prb\{|Z-m|\le\rho t\}
 \ge\frac{2\rho}{e\sqrt{2\pi}}e^{-m^2/t^2}.
\]
This follows by bounding the Gaussian density throughout the interval,
using $(|m|/t+\rho)^2/2\le m^2/t^2+\rho^2$.
Multiply the interval estimates and the bulk probability $1/2$.
Since $N/t^2\ge1$ and $t\le1$, the result is at least
$\exp(-ND_*/t^2)$ with \eqref{eq:target-constants}.
\end{proof}


\end{document}
